\documentclass[12pt,letterpaper]{article}
\usepackage[spanish,es-nodecimaldot]{babel}
\usepackage[utf8]{inputenc}
\usepackage[T1]{fontenc}
\usepackage{lmodern}
\usepackage{amsmath,amssymb}
\usepackage{geometry}
\usepackage{enumitem}
\usepackage{tcolorbox}
\usepackage{xcolor}
\geometry{margin=2.5cm}
\setlength{\parskip}{0.5em}
\setlength{\parindent}{0pt}

\title{Solución exacta paso a paso de la ecuación diferencial del modelo}
\author{}
\date{}

\begin{document}
\maketitle

\section*{1. Planteamiento del modelo}

A partir de la segunda ley de Newton para un cuerpo que cae bajo la acción de la gravedad y una fuerza de resistencia proporcional a la velocidad, se obtiene:

\[
\frac{dv}{dt}=\frac{mg-cv}{m}
\]

Dividiendo cada término entre la masa \(m\):

\[
\frac{dv}{dt}=g-\frac{c}{m}v
\]

Esta es una ecuación diferencial ordinaria de primer orden para la velocidad \(v(t)\):

\[
\boxed{\frac{dv}{dt}=g-\frac{c}{m}v}
\]

\begin{tcolorbox}[colback=blue!4,colframe=blue!60!black,title=Datos del modelo]
\begin{itemize}[leftmargin=*]
    \item Variable dependiente: \(v(t)\), velocidad del cuerpo.
    \item Variable independiente: \(t\), tiempo.
    \item Parámetros: \(m\), masa del cuerpo; \(c\), coeficiente de resistencia.
    \item Aceleración gravitacional: \(g\).
\end{itemize}
\end{tcolorbox}

\section*{2. Organización de la ecuación diferencial}

La ecuación se escribe en forma lineal:

\[
\frac{dv}{dt}+\frac{c}{m}v=g
\]

Esta ecuación tiene la forma general:

\[
\frac{dy}{dx}+P(x)y=Q(x)
\]

Comparando:

\[
P(t)=\frac{c}{m}, \qquad Q(t)=g
\]

Por tanto, puede resolverse mediante el método del factor integrante.

\section*{3. Cálculo del factor integrante}

El factor integrante se define como:

\[
\mu(t)=e^{\int P(t)\,dt}
\]

Como \(P(t)=\frac{c}{m}\), entonces:

\[
\mu(t)=e^{\int \frac{c}{m}\,dt}
\]

\[
\mu(t)=e^{\frac{c}{m}t}
\]

Por lo tanto:

\[
\boxed{\mu(t)=e^{\frac{c}{m}t}}
\]

\section*{4. Multiplicación por el factor integrante}

Multiplicamos toda la ecuación diferencial por \(e^{\frac{c}{m}t}\):

\[
e^{\frac{c}{m}t}\frac{dv}{dt}+\frac{c}{m}e^{\frac{c}{m}t}v=ge^{\frac{c}{m}t}
\]

El lado izquierdo corresponde a la derivada del producto:

\[
\frac{d}{dt}\left(v e^{\frac{c}{m}t}\right)=e^{\frac{c}{m}t}\frac{dv}{dt}+\frac{c}{m}e^{\frac{c}{m}t}v
\]

Por tanto, la ecuación queda:

\[
\frac{d}{dt}\left(v e^{\frac{c}{m}t}\right)=ge^{\frac{c}{m}t}
\]

\section*{5. Integración de ambos lados}

Integramos con respecto al tiempo:

\[
\int \frac{d}{dt}\left(v e^{\frac{c}{m}t}\right)dt=\int ge^{\frac{c}{m}t}dt
\]

Entonces:

\[
v e^{\frac{c}{m}t}=g\int e^{\frac{c}{m}t}dt
\]

Como:

\[
\int e^{\frac{c}{m}t}dt=\frac{m}{c}e^{\frac{c}{m}t}
\]

se obtiene:

\[
v e^{\frac{c}{m}t}=\frac{mg}{c}e^{\frac{c}{m}t}+C
\]

\section*{6. Despeje de la velocidad}

Dividimos toda la expresión entre \(e^{\frac{c}{m}t}\):

\[
v(t)=\frac{mg}{c}+Ce^{-\frac{c}{m}t}
\]

Por tanto, la solución general exacta es:

\[
\boxed{v(t)=\frac{mg}{c}+Ce^{-\frac{c}{m}t}}
\]

\section*{7. Aplicación de una condición inicial}

Si se conoce la velocidad inicial del cuerpo:

\[
v(0)=v_0
\]

entonces se sustituye \(t=0\) en la solución general:

\[
v(0)=\frac{mg}{c}+Ce^{-\frac{c}{m}(0)}
\]

Como \(e^0=1\):

\[
v_0=\frac{mg}{c}+C
\]

Despejando \(C\):

\[
C=v_0-\frac{mg}{c}
\]

Por tanto, la solución exacta con condición inicial \(v(0)=v_0\) es:

\[
\boxed{v(t)=\frac{mg}{c}+\left(v_0-\frac{mg}{c}\right)e^{-\frac{c}{m}t}}
\]

\section*{8. Caso particular: el cuerpo parte del reposo}

Si el cuerpo se suelta desde el reposo, entonces:

\[
v(0)=0
\]

Por tanto:

\[
0=\frac{mg}{c}+C
\]

\[
C=-\frac{mg}{c}
\]

Sustituyendo en la solución general:

\[
v(t)=\frac{mg}{c}-\frac{mg}{c}e^{-\frac{c}{m}t}
\]

Factorizando:

\[
\boxed{v(t)=\frac{mg}{c}\left(1-e^{-\frac{c}{m}t}\right)}
\]

Esta es la solución exacta para un cuerpo que cae desde el reposo con resistencia lineal del aire.

\section*{9. Interpretación física}

Cuando \(t\to\infty\), se cumple que:

\[
e^{-\frac{c}{m}t}\to 0
\]

Entonces:

\[
\lim_{t\to\infty}v(t)=\frac{mg}{c}
\]

Por consiguiente, la velocidad terminal del cuerpo es:

\[
\boxed{v_T=\frac{mg}{c}}
\]

Esto significa que, después de cierto tiempo, la velocidad deja de aumentar significativamente y se aproxima al valor constante \(\frac{mg}{c}\).

\section*{10. Verificación de la solución}

Tomemos la solución para el caso en que el cuerpo parte del reposo:

\[
v(t)=\frac{mg}{c}\left(1-e^{-\frac{c}{m}t}\right)
\]

Derivando:

\[
\frac{dv}{dt}=\frac{mg}{c}\left(\frac{c}{m}e^{-\frac{c}{m}t}\right)
\]

\[
\frac{dv}{dt}=ge^{-\frac{c}{m}t}
\]

Ahora calculamos el lado derecho de la ecuación diferencial original:

\[
g-\frac{c}{m}v
\]

Sustituyendo \(v(t)\):

\[
g-\frac{c}{m}\left[\frac{mg}{c}\left(1-e^{-\frac{c}{m}t}\right)\right]
\]

Simplificando:

\[
g-g\left(1-e^{-\frac{c}{m}t}\right)
\]

\[
g-g+ge^{-\frac{c}{m}t}
\]

\[
ge^{-\frac{c}{m}t}
\]

Como:

\[
\frac{dv}{dt}=ge^{-\frac{c}{m}t}
\]

y:

\[
g-\frac{c}{m}v=ge^{-\frac{c}{m}t}
\]

la solución satisface la ecuación diferencial. Por tanto, queda verificada.

\section*{11. Resultado final}

La solución exacta general es:

\[
\boxed{v(t)=\frac{mg}{c}+\left(v_0-\frac{mg}{c}\right)e^{-\frac{c}{m}t}}
\]

Si el cuerpo parte del reposo, la solución exacta es:

\[
\boxed{v(t)=\frac{mg}{c}\left(1-e^{-\frac{c}{m}t}\right)}
\]

Y la velocidad terminal es:

\[
\boxed{v_T=\frac{mg}{c}}
\]

\end{document}
